\begin{table}[htbp]
\centering
\caption{Médias de rendimento por fonte nos municípios com maiores e menores rendimentos totais, Sergipe, 2022}
\label{tab:item4-renda-fontes}
\scriptsize
\begin{tabular}{lrrrrr}
\toprule
Município & Trabalho (R\$) & Transferências (R\$) & Total (R\$) & Trabalho (\%) & Transferências (\%) \\
\midrule
\multicolumn{6}{l}{\textbf{Maiores médias de rendimento total}} \\
Aracaju & 3.371,02 & 205,71 & 3.576,73 & 94,25 & 5,75 \\
Barra dos Coqueiros & 2.784,52 & 152,94 & 2.937,46 & 94,79 & 5,21 \\
São Cristóvão & 1.998,48 & 117,16 & 2.115,64 & 94,46 & 5,54 \\
\midrule
\multicolumn{6}{l}{\textbf{Menores médias de rendimento total}} \\
Ilha das Flores & 975,29 & 99,38 & 1.074,66 & 90,75 & 9,25 \\
Santa Luzia do Itanhy & 957,07 & 113,26 & 1.070,34 & 89,42 & 10,58 \\
Tomar do Geru & 858,82 & 130,92 & 989,74 & 86,77 & 13,23 \\
\bottomrule
\end{tabular}
\fonte{Elaboração própria com base nos microdados do Censo Demográfico 2022.}
\nota{As médias e participações são as estimadas pelo código para os registros válidos conjuntamente nas três variáveis de rendimento.}
\end{table}
